% ============================================================
% SECCION 5 - CONCLUSIONES (SINTETIZADA)
% ============================================================

\hypertarget{sec5}{\section{Conclusiones}}\label{sec5}

\subsection{Síntesis de Hallazgos}

\begin{keypoint}
\textbf{Cinco conclusiones críticas:}
\begin{enumerate}
    \itemsep=0.15cm
    \item \textbf{Escala}: 72,9\% encontró desinformación; sistemática, coordinada y financiada
    \item \textbf{Multinivel}: Desde individual hasta corporaciones transnacionales
    \item \textbf{Arma política}: Dirigida a desfavorecer procesos progresistas
    \item \textbf{Amenaza democrática}: 80\% latinoamericanos la ve como riesgo
    \item \textbf{Vacío legal}: Chile sin marco específico (UE y Brasil sí)
\end{enumerate}
\end{keypoint}

\subsection{Implicaciones}

\begin{itemize}
    \itemsep=0.12cm
    \item \textbf{Erosión de confianza}: Debilita instituciones y proceso electoral
    \item \textbf{Polarización}: Burbujas informativas; diálogo imposible
    \item \textbf{Tribalismo}: Se acepta acusación sin verificación
    \item \textbf{Ataque prensa}: 76\% ataques a mujeres periodistas
\end{itemize}

\subsection{Estrategias de Defensa}

\subsubsection{Individual}
Educación preventiva, recordatorios de precisión, escepticismo informado

\subsubsection{Institucional}
UE: multas 10\%. Brasil: responsabilidad proactiva. Chile: urgencia legislativa

\subsubsection{Colaborativo}
Red de 40+ medios, transparencia corporativa, academia participativa

\subsection{Desafíos}

\textbf{Paradoja:} Personas educadas justifican mejor la desinformación que les acomoda.

\textbf{Legitimidad:} ¿Quién regula? ¿Estado, plataformas, sociedad civil?

\textbf{Tecnología:} IA crea desinformación masiva; respuestas sofisticadas necesarias

\subsection{Recomendaciones}

\begin{description}
    \item[Política:] Legislación específica (modelo europeo adaptado). Autoridad fiscalizadora independiente. Transparencia de financiamiento digital.
    
    \item[Plataformas:] Transparencia algorítmica. Responsabilidad genuina por contenido. Colaboración con verificadores. Desmonetización de desinformación.
    
    \item[Sociedad Civil:] Media literacy en educación formal. Cuestionamiento de información emocional. Valentía de desacreditar falsedad propia. Fact-checking participativo.
    
    \item[Academia:] Investigación continua. Detección automática de bots. Participación activa en debate público. Investigación interdisciplinaria.
\end{description}

\subsection{Reflexión Final}

Desinformación: problema \textbf{tecnológico en forma, humano en fondo}. Se apropia de divisiones, sesgos y miedos, amplificándolos digitalmente.

Sin solución única. Va desde educar al niño sobre media literacy hasta regular gigantes tecnológicos. En el medio: balance entre libertad de expresión y seguridad informativa.

\textbf{La inacción no es opción.} Cuando 80\% ve amenaza, corporaciones manipulan electorales, bots distorsionan realidad: \textbf{obligación actuar}.

\textbf{Pregunta: no si regulamos, sino cuándo, cómo, y con qué valores.}
